\section{Introduction}
\IEEEPARstart{V}{ideo} violence classification asks whether a short sequence contains violent activity. The problem is difficult when several people overlap, when a camera moves, or when nonviolent activity contains rapid motion and close physical proximity. A cheering crowd can produce strong image motion, while an aggressive interaction can occupy only a small part of a frame. A clip classifier must therefore separate the label-relevant signal from recording conditions, detection behavior, and unrelated activity.

Motion descriptors and learned video representations provide established approaches to this problem \cite{hassner2012,nievas2011,serrano2015,cheng2020}. Human-pose representations offer a complementary description in terms of joints and trajectories. They can suppress some appearance information, but they introduce estimation errors, missing observations, and association decisions. A confidence channel may help distinguish reliable from unreliable joints, yet the same channel can also expose properties of the scene or recording domain. Relative-person geometry may describe interaction, but perspective and tracking errors can distort that geometry. These additions require evaluation rather than an assumed benefit.

The primary question of this study is whether explicit pose reliability and nearest-person interactions improve a fixed temporal baseline on crowd videos. We implement four feature variants using the same pose caches, source assignments, encoder family, and training budget. We also examine a correction to body centering, which otherwise changes abruptly when hip joints become unavailable. Matched fits compare the original fallback with a continuous offset estimate. The correction addresses feature construction and is distinct from a claim of improved external generalization.

A second question concerns the information that makes these classifiers effective. A seven-variable detection-quality classifier, a motion reference, a coverage ablation, and matched input controls test different aspects of the available signal. In particular, zeroing kinematic channels leaves confidence, validity, observation masks, neighbor counts, and coverage accessible. Shuffling derived slices disrupts their order but retains velocities already calculated from adjacent observations. These distinctions determine what conclusions the controls can support.

The study then adds a frozen pretrained video reference on the same crowd folds and a separate supervised experiment on RWF-2000. The latter fits a new classifier using the surveillance training release and develops one limited fine-tuning candidate. Validation alone selects its checkpoint and threshold. A fixed external Hockey Fight experiment compares that candidate, the original frozen surveillance reference, and a crowd-trained pose checkpoint. This follow-up tests recording-domain transfer and the consequences of changing the operating point, while preserving the original pose question.

The contributions are empirical and implementation-focused. First, we document a reproducible multi-person pose pipeline with confidence and nearest-person variants and a matched correction of the coordinate origin. Second, we report controlled evidence that detection-related inputs remain predictive and explain why an ordering control does not remove motion. Third, we provide separate crowd and surveillance video references and a validation-selected fine-tuning comparison with explicit miss and false-alarm counts. Finally, we connect aggregate results to decoded diagnostic videos, runtime measurements, and saved evaluation records.

The resulting evidence is mixed. Interaction features give the highest mean accuracy among the four corrected pose variants, but the combined variant performs worse than the baseline. The quality-only control is competitive, and frozen video features give a stronger crowd reference. Fine-tuning improves Fight recall on the surveillance evaluation, while slightly reducing accuracy and increasing false alarms. External pose transfer is poor. These findings describe the implemented models and experimental conditions, without establishing a novel general-purpose architecture or a deployment-ready detector.

The manuscript treats the evaluations as exploratory. Existing crowd and RWF footage had been inspected during earlier development, known source links do not exhaust all possible overlap, and independent label review remains pending. We retain those qualifications throughout the methods and results. The task is whole-clip binary classification. The outputs do not estimate the exact onset of violence, attribute responsibility to a person, or validate performance on an uninterrupted surveillance stream.
